% ==============================================================================
% KHOẢN KHUNG THỰC NGHIỆM ĐO ĐẠC TỰ ĐỘNG TỪ BENCHMARK.PY (CRAM-ABSA)
% Định dạng rút gọn (Compact Anti-Overflow) chuẩn Springer Nature
% ==============================================================================

% TABLE 1 (RQ1 COMPACT): ORDINAL STAR CLASSIFICATION (6 COLUMNS - ANTI-OVERFLOW)
\begin{table}[htbp]
\centering
\caption{Khuôn khổ đối chứng phân loại thứ bậc 5 mức sao (RQ1) trên tập kiểm thử.}
\label{tab:benchmark_rq1}
\small
\begin{tabular}{lccccc}
\hline
\textbf{Phương pháp} & \textbf{Acc@0 (\%)} & \textbf{Acc@1 (\%)} & \textbf{Macro-MAE} & \textbf{QWK} & \textbf{Macro-F1 (\%)} \\
\hline
Multi-Class CE Baseline & -- & -- & -- & -- & -- \\
Standard Ordinal Reg.   & -- & -- & -- & -- & -- \\
\textbf{CRAM-ABSA (Proposed)} & \textbf{--} & \textbf{--} & \textbf{--} & \textbf{--} & \textbf{--} \\
\hline
\end{tabular}

\vspace{2pt}
\raggedright \footnotesize \textit{Ghi chú:} Macro-MAE tính trung bình sai số tuyệt đối trên từng mức sao 1$\star \to 5\star$ nhằm kiểm chứng độ lệch trên các lớp thiểu số mà không làm tràn cột bảng.
\end{table}

% TABLE 2 (RQ2 COMPACT): CROSS-ASPECT LEAKAGE & RATIONALE TOKEN F1 (4 COLUMNS)
\begin{table}[htbp]
\centering
\caption{Khuôn khổ khảo sát triệt tiêu rò rỉ căn cứ liên khía cạnh và chất lượng rationale (RQ2).}
\label{tab:benchmark_rq2}
\small
\begin{tabular}{lccc}
\hline
\textbf{Khía cạnh (Aspect)} & \textbf{Token F1 (Gốc)} & \textbf{Token F1 (+$\mathcal{L}_{\text{overlap}}$)} & \textbf{Tỷ lệ Rò rỉ (\%)} \\
\hline
Service    & -- & \textbf{--} & \textbf{--} \\
Facility   & -- & \textbf{--} & \textbf{--} \\
Experience & -- & \textbf{--} & \textbf{--} \\
Amenity    & -- & \textbf{--} & \textbf{--} \\
Loyalty    & -- & \textbf{--} & \textbf{--} \\
Branding   & -- & \textbf{--} & \textbf{--} \\
\hline
\textbf{Trung bình (Macro-Avg)} & -- & \textbf{--} & \textbf{--} \\
\hline
\end{tabular}

\vspace{2pt}
\raggedright \footnotesize \textit{Ghi chú:} Tỷ lệ rò rỉ căn cứ phản ánh mức độ chồng lấn mặt nạ token ngoài ý muốn giữa các khía cạnh phân tách.
\end{table}

% TABLE 3 (RQ3 COMPACT): GENERATIVE QA ON 60 VERIFIED QUESTIONS (6 COLUMNS)
\begin{table*}[htbp]
\centering
\caption{Khuôn khổ đánh giá chất lượng hội thoại và độ trễ suy luận (RQ3) trên bộ 60 câu hỏi kiểm định.}
\label{tab:benchmark_rq3}
\small
\begin{tabular}{lccccc}
\hline
\textbf{Cấu hình Phương pháp} & \textbf{ROUGE-1 (\%)} & \textbf{ROUGE-2 (\%)} & \textbf{ROUGE-L (\%)} & \textbf{Token F1 (\%)} & \textbf{Độ trễ (ms)} \\
\hline
Vanilla BART-base (No RAG)         & -- & -- & -- & -- & -- \\
Standard Text RAG (No Time-Weight) & -- & -- & -- & -- & -- \\
Metadata Recency RAG (Zero-Shot)   & -- & -- & -- & -- & -- \\
\textbf{CRAM-ABSA Hybrid (Proposed)} & \textbf{--} & \textbf{--} & \textbf{--} & \textbf{--} & \textbf{--} \\
\hline
\end{tabular}

\vspace{2pt}
\raggedright \footnotesize \textit{Ghi chú:} Bộ 60 câu hỏi bao phủ đồng đều 6 khía cạnh dịch vụ tại 35 khách sạn thực tế tại Việt Nam.
\end{table*}

% MASTER SUMMARY TABLE (SIÊU RÚT GỌN - DỄ DÀNG CHÈN MỌI VỊ TRÍ)
\begin{table}[htbp]
\centering
\caption{Bảng tổng hợp chỉ số cốt lõi của mô hình đề xuất CRAM-ABSA qua cả 3 câu hỏi nghiên cứu.}
\label{tab:benchmark_master_summary}
\small
\begin{tabular}{lcccccc}
\hline
\textbf{Mô hình} & \textbf{Acc@0} & \textbf{Macro-MAE} & \textbf{QWK} & \textbf{ROUGE-L} & \textbf{Token F1} & \textbf{Độ trễ (ms)} \\
\hline
\textbf{CRAM-ABSA (Proposed)} & --\% & -- & -- & --\% & --\% & -- \\
\hline
\end{tabular}
\end{table}
